%%%%%%%%%%%%%%%%%
% This is an sample CV template created using altacv.cls
% (v1.7.4, 30 July 2025) written by LianTze Lim (liantze@gmail.com). Compiles with pdfLaTeX, XeLaTeX and LuaLaTeX.
%
%% It may be distributed and/or modified under the
%% conditions of the LaTeX Project Public License, either version 1.3
%% of this license or (at your option) any later version.
%% The latest version of this license is in
%%    http://www.latex-project.org/lppl.txt
%% and version 1.3 or later is part of all distributions of LaTeX
%% version 2003/12/01 or later.
%%%%%%%%%%%%%%%%

%% Use the "normalphoto" option if you want a normal photo instead of cropped to a circle
% \documentclass[10pt,a4paper,withhyper,normalphoto]{altacv}

\documentclass[10pt,a4paper,withhyper]{altacv}
%% AltaCV uses the fontawesome5 and simpleicons packages.
%% See http://texdoc.net/pkg/fontawesome5 and http://texdoc.net/pkg/simpleicons for full list of symbols.

% Change the page layout if you need to
\geometry{left=1.25cm,right=1.25cm,top=1.5cm,bottom=1.5cm,columnsep=1.2cm}

% The paracol package lets you typeset columns of text in parallel
\usepackage{paracol}
\usepackage{ragged2e}
\usepackage{xcolor}
\definecolor{maroon}{RGB}{128, 0, 0}

\sloppy
\emergencystretch=2em

% Change the font if you want to, depending on whether
% you're using pdflatex or xelatex/lualatex
% WHEN COMPILING WITH XELATEX PLEASE USE
% xelatex -shell-escape -output-driver="xdvipdfmx -z 0" sample.tex
\iftutex
  % If using xelatex or lualatex:
  \setmainfont{Roboto Slab}
  \setsansfont{Lato}
  \renewcommand{\familydefault}{\sfdefault}
\else
  % If using pdflatex:
  \usepackage[rm]{roboto}
  \usepackage[defaultsans]{lato}
  % \usepackage{sourcesanspro}
  \renewcommand{\familydefault}{\sfdefault}
\fi

% Change the colours if you want to
\definecolor{SlateGrey}{HTML}{2E2E2E}
\definecolor{LightGrey}{HTML}{666666}
\definecolor{DarkPastelRed}{HTML}{450808}
\definecolor{PastelRed}{HTML}{8F0D0D}
\definecolor{GoldenEarth}{HTML}{E7D192}
\colorlet{name}{black}
\colorlet{tagline}{PastelRed}
\colorlet{heading}{DarkPastelRed}
\colorlet{headingrule}{GoldenEarth}
\colorlet{subheading}{PastelRed}
\colorlet{accent}{PastelRed}
\colorlet{emphasis}{SlateGrey}
\colorlet{body}{LightGrey}

% Change some fonts, if necessary
\renewcommand{\namefont}{\Huge\rmfamily\bfseries}
\renewcommand{\personalinfofont}{\footnotesize}
\renewcommand{\cvsectionfont}{\LARGE\rmfamily\bfseries}
\renewcommand{\cvsubsectionfont}{\large\bfseries}


% Change the bullets for itemize and rating marker
% for \cvskill if you want to
\renewcommand{\cvItemMarker}{{\small\textbullet}}
\renewcommand{\cvRatingMarker}{\faCircle}
% ...and the markers for the date/location for \cvevent
% \renewcommand{\cvDateMarker}{\faCalendar*[regular]}
% \renewcommand{\cvLocationMarker}{\faMapMarker*}


% If your CV/résumé is in a language other than English,
% then you probably want to change these so that when you
% copy-paste from the PDF or run pdftotext, the location
% and date marker icons for \cvevent will paste as correct
% translations. For example Spanish:
% \renewcommand{\locationname}{Ubicación}
% \renewcommand{\datename}{Fecha}


%% Use (and optionally edit if necessary) this .tex if you
%% want to use an author-year reference style like APA(6)
%% for your publication list
% \input{pubs-authoryear}

%% Use (and optionally edit if necessary) this .tex if you
%% want an originally numerical reference style like IEEE
%% for your publication list
\input{pubs-num}

%% sample.bib contains your publications
\addbibresource{sample.bib}

\begin{document}
\name{Faraz Islam}
\tagline{\small Experienced Cloud Platform Engineer @ Veeam Software \textcolor{black}{| IT Cloud Systems \& Infrastructure Specialist | Azure DevOps | Azure Cloud Management | AWS | GitHub | IaC | Cloud Monitoring and Security}
}

%% You can add multiple photos on the left or right
\photoR{2.8cm}{FarazCV}
% \photoL{2.5cm}{Yacht_High,Suitcase_High}

\personalinfo{%
  % Not all of these are required!
  \email{fislam2010@gmail.com}
  \phone{+4917620705482}
  \mailaddress{Eschersheimer Landstr. 461, 60431 Frankfurt am Main}
  \location{Germany}
  \linkedin{https://www.linkedin.com/in/faraz-islam/}
  \homepage{https://islamfaraz.github.io/portfolio/}
  %% You can add your own arbitrary detail with
  %% \printinfo{symbol}{detail}[optional hyperlink prefix]
  % \printinfo{\faPaw}{Hey ho!}[https://example.com/]

  %% Or you can declare your own field with
  %% \NewInfoFiled{fieldname}{symbol}[optional hyperlink prefix] and use it:
  %%
  %% For services and platforms like Mastodon where there isn't a
  %% straightforward relation between the user ID/nickname and the hyperlink,
  %% you can use \printinfo directly e.g.
  % \printinfo{\faMastodon}{@username@instace}[https://instance.url/@username]
  %% But if you absolutely want to create new dedicated info fields for
  %% such platforms, then use \NewInfoField* with a star:
}

\makecvheader
%% Depending on your tastes, you may want to make fonts of itemize environments slightly smaller
% \AtBeginEnvironment{itemize}{\small}

%% Set the left/right column width ratio to 6:4.
\columnratio{0.6}

% Start a 2-column paracol. Both the left and right columns will automatically
% break across pages if things get too long.
\begin{paracol}{2}
\cvsection{Experience}

\cvevent{Experienced Cloud Platform Engineer}{Veeam Software}{Aug 2024 -- Present}{Germany}
\begin{justify}
\begin{itemize}

    \item Developed Azure Landing Zone–aligned architectures using Terraform and Bicep, increasing provisioning consistency by 70\% and cutting environment setup time from days to under 2 hours.
    \item Standardized AKS clusters across environments, decreasing configuration drift by 60\% and boosting deployment reliability by 40\%.
     \item Implemented AWS EKS clusters with secure networking (VPC, PrivateLink, Security Groups), improving workload isolation and achieving 80\% compliance.
    \item Designed hybrid Azure–AWS architectures for cross-cloud workloads, improving service resilience by 30\% and reducing vendor lock-in.
    \item Facilitated architecture reviews and technical workshops, accelerating cloud adoption and lowering design ambiguity by 50\%.
    \item Created umbrella CI/CD pipelines in Azure DevOps for coordinated deployment of 20+ microservices, decreasing release effort by 50\% and strengthening traceability.
    \item Established OpenTelemetry collectors and Grafana dashboards, refining MTTR by 35\% and enabling real-time observability.
    \item Optimized Docker images and artifact workflows, lessening image size by up to 45\% and enhancing build times by 30\%.
    \item Enforced platform security using NSGs, Private DNS Zones, and Azure Policies, achieving 75\% policy compliance.
    \item Systematized operational notifications via Power Automate, boosting stakeholder visibility and lowering manual communication overhead by 80\%.
    \item Built Terraform modules supporting both Azure and AWS, standardizing provisioning across clouds and reducing IaC duplication by 40\%.
    \item Integrated CloudWatch and AWS X-Ray with OpenTelemetry pipelines, improving distributed tracing coverage by 45\%.
    \item Optimized AWS compute and storage costs using rightsizing and lifecycle policies, reducing monthly spend by 20\%.
    
\end{itemize}
\end{justify}
\divider

\cvevent{Deputy Head of Digital Experience Department}{Quipu GmbH}{Apr 2024 -- Jul 2024}{Germany}
\begin{justify}
\begin{itemize}

    \item Directed cloud modernization initiatives using Azure, elevating platform reliability and performance for 10+ digital services.
    \item Mentored cross-functional engineering teams, increasing delivery velocity by 25\% through refined DevOps practices.
    \item Advanced adoption of Azure DevOps CI/CD pipelines, decreasing deployment friction and enhancing release frequency by 40\%.
    \item Collaborated with business stakeholders to translate requirements into scalable Azure architectures, lowering rework by 30\%.
    \item Introduced governance frameworks that boosted service uptime and decreased incident volume by 20\%.
    
\end{itemize}
\end{justify}
\divider

\cvevent{Coordinator of Automations}{Quipu GmbH}{Nov 2022 -- Mar 2024}{Germany}
\begin{justify}
\begin{itemize}

    \item Led a globally distributed automation engineering team, enhancing collaboration and delivery efficiency by 30\%.
    \item Applied Infrastructure as Code (IaC) using Terraform and Azure DevOps, decreasing provisioning time from hours to minutes.
    \item Modernized hybrid cloud environments, lowering operational incidents by 25\% and strengthening system stability.
    \item Refined CI/CD workflows with Azure DevOps Boards and Pipelines, boosting deployment success rates by 35\%.
    \item Delivered automation solutions that saved 50+ hours per month across teams.
    
\end{itemize}
\end{justify}
\divider

\cvevent{Agile Automation Engineer}{Quipu GmbH}{May 2018 -- Oct 2022}{Germany}
\begin{justify}
\begin{itemize}

    \item Engineered automation solutions across Azure PaaS, SaaS, and IaaS, enhancing operational efficiency by 30\%.
    \item Developed microservices and cloud-native applications, lowering service latency and improving scalability.
    \item Executed CI/CD pipelines that decreased deployment time by 40\% and strengthened release reliability.
    \item Supported Azure Virtual Desktop, API Management, and enterprise cloud workloads, boosting service availability and user satisfaction.
    
\end{itemize}
\end{justify}
\divider

\cvevent{Student Intern}{Quipu GmbH}{May 2017 -- May 2018}{Germany}
\begin{justify}
\begin{itemize}

    \item Performed statistical data analysis and cloud administration tasks, refining reporting accuracy by 20\%.
    \item Assisted Azure ML experiments and data modeling workflows, accelerating research cycles.
    \item Enhanced ServiceDesk operations by streamlining ticket triage and decreasing resolution time by 15\%
    
\end{itemize}
\end{justify}

\divider


\cvsection{Projects}

\cvevent{\textbf{\textcolor{maroon}{Azure Landing Zone Modernization Program}}}{}{}{}
\begin{justify}
\begin{itemize}
    \item Modernized the enterprise Azure foundation by introducing modular Landing Zone components, reducing onboarding time for new workloads by 60\% and improving architectural consistency across teams.
    \item Implemented a governance and identity model that increased policy compliance to 85\% and reduced audit findings by 40\%.
    \item Delivered reusable IaC patterns that cut environment provisioning effort by 50\% and improved platform scalability for future expansions.
\end{itemize}
\end{justify}

\divider

\cvevent{\textbf{\textcolor{maroon}{AKS Platform Reliability \& Observability Initiative}}}{}{}{}
\begin{justify}
\begin{itemize}
    \item Designed a reliability framework for AKS workloads, improving service resilience by 35\% through standardized deployment patterns, health probes, and scaling strategies.
    \item Built an observability blueprint integrating OpenTelemetry and metrics pipelines, reducing MTTR by 35\% and enabling proactive detection of performance issues.
    \item Partnered with development teams to optimize microservice performance, lowering latency by up to 25\% and improving user experience across critical applications.

\end{itemize}
\end{justify}
\medskip

\cvevent{\textbf{\textcolor{maroon}{Enterprise CI/CD Standardization Program}}}{}{}{}
\begin{justify}
\begin{itemize}
    \item Defined a unified CI/CD architecture for microservices, increasing release reliability by 30\% and reducing manual deployment effort by 50\%.
    \item Introduced artifact governance and automated quality gates, decreasing rollback frequency by 45\% and improving traceability across releases.
    \item Established cross-team deployment orchestration patterns that reduced operational coordination overhead by 40\% and improved delivery predictability.
\end{itemize}
\end{justify}

\divider

\cvsection{A Day of My Life}

% Adapted from @Jake's answer from http://tex.stackexchange.com/a/82729/226
% \wheelchart{outer radius}{inner radius}{
% comma-separated list of value/text width/color/detail}
\wheelchart{1.5cm}{0.5cm}{%
  6/8em/accent!30/{Sleep,\\beautiful sleep},
  3/8em/accent!40/Hopeful novelist by night,
  8/8em/accent!60/Actively learning \& upskilling,
  2/10em/accent/Sports \& relaxation,
  5/6em/accent!20/Spending time with family
}

% use ONLY \newpage if you want to force a page break for
% ONLY the current column
\newpage

% \cvsection{Publications}

% %% Specify your last name(s) and first name(s) as given in the .bib to automatically bold your own name in the publications list.
% %% One caveat: You need to write \bibnamedelima where there's a space in your name for this to work properly; or write \bibnamedelimi if you use initials in the .bib
% %% You can specify multiple names, especially if you have changed your name or if you need to highlight multiple authors.
% \mynames{Lim/Lian\bibnamedelima Tze,
%   Wong/Lian\bibnamedelima Tze,
%   Lim/Tracy,
%   Lim/L.\bibnamedelimi T.}
% %% MAKE SURE THERE IS NO SPACE AFTER THE FINAL NAME IN YOUR \mynames LIST

% \nocite{*}

% \printbibliography[heading=pubtype,title={\printinfo{\faBook}{Books}},type=book]

% \divider

% \printbibliography[heading=pubtype,title={\printinfo{\faFile*[regular]}{Journal Articles}},type=article]

% \divider

% \printbibliography[heading=pubtype,title={\printinfo{\faUsers}{Conference Proceedings}},type=inproceedings]

%% Switch to the right column. This will now automatically move to the second
%% page if the content is too long.
\switchcolumn

\cvsection{My Life Philosophy}

\begin{quote}
``Every day is an opportunity to learn, improve, and help others succeed.''
\end{quote}

% \cvsection{Most Proud of}

% \cvachievement{\faTrophy}{Fantastic Achievement}{and some details about it}

% \divider

% \cvachievement{\faHeartbeat}{Another achievement}{more details about it of course}

% \divider

% \cvachievement{\faHeartbeat}{Another achievement}{more details about it of course}

\cvsection{Education}

\cvevent{Ph.D.\ in Software Applications}{Integral University, India}{January 2023 -- Present}{}

\divider

\cvevent{M.Sc.\ in High Integrity Systems}{Frankfurt University of Applied Sciences, Germany}{April 2016 -- December 2020}{}

\divider

\cvevent{B.Tech.\ in Computer Science}{Aligarh Muslim University, India}{July 2011 -- June 2015}{}


\cvsection{SKILLS}

% Don't overuse these \cvtag boxes — they're just eye-candies and not essential. If something doesn't fit on a single line, it probably works better as part of an itemized list (probably inlined itemized list), or just as a comma-separated list of strengths.

\textbf{\textcolor{maroon}{Cloud Architecture:}}\\[4pt]
\cvtag{Azure} 
\cvtag{AWS} 
\cvtag{Landing Zones} 
\cvtag{Well-Architecture Framework} 
\cvtag{Microservices} 
\cvtag{Event-Driven Design}

\divider
\textbf{\textcolor{maroon}{Containers \& Orchestration:}}\\ [4pt]
\cvtag{AKS}
\cvtag{EKS}
\cvtag{Docker}
\cvtag{Helm}
\cvtag{Service Mesh}

\divider
\textbf{\textcolor{maroon}{Infrastructure as Code}}\\ [4pt]
\cvtag{Terraform}
\cvtag{Bicep}
\cvtag{CloudFormation}

\divider
\textbf{\textcolor{maroon}{DevOps \& CI/CD:}}\\
\cvtag{Azure DevOps}
\cvtag{GitHub Actions}
\cvtag{TFS}
\cvtag{CI/CD}
\cvtag{YAML Pipelines}
\cvtag{JFrog Artifactory}
\cvtag{ADO Dashboards}

\divider
\textbf{\textcolor{maroon}{Automation \& Scripting:}}\\
\cvtag{PowerShell}
\cvtag{Bash}
\cvtag{Power Automate}
\cvtag{Azure CLI}
\cvtag{AWS CLI}
\cvtag{Azure PowerShell}

\divider
\textbf{\textcolor{maroon}{Networking \& Security:}}\\
\cvtag{VPC / VNet}
\cvtag{PrivateLink}
\cvtag{Private Link Service}
\cvtag{Security Groups}
\cvtag{NSG}
\cvtag{Az Firewall}
\cvtag{AWS WAF}
\cvtag{Azure WAF}

\divider
\textbf{\textcolor{maroon}{Observability:}}\\
\cvtag{SolarWinds}
\cvtag{OpenTelemetry}
\cvtag{Prometheus}
\cvtag{Grafana}
\cvtag{Elastic Search}
\cvtag{Application Insights}
\cvtag{Azure Monitor}
\cvtag{CloudWatch}
\cvtag{Log Analytics}

\divider
\textbf{\textcolor{maroon}{Cost Optimization:}}\\
\cvtag{Azure Cost Management}
\cvtag{AWS Cost Explorer}
\cvtag{Rightsizing}
\cvtag{Reserved Instances}
\cvtag{Cloud Cost Planning}

\divider
\textbf{\textcolor{maroon}{Identity \& Access Management:}}\\
\cvtag{Entra ID}
\cvtag{Conditional Access}
\cvtag{AWS IAM}
\cvtag{Azure RBAC}
\cvtag{AWS SSO}
\cvtag{Identity Governance}

\divider
\textbf{\textcolor{maroon}{Governance \& Compliance:}}\\
\cvtag{Policy as Code}
\cvtag{CIS Benchmarks}
\cvtag{Zero Trust}
\cvtag{Azure Policies}
\cvtag{AWS Config}

\divider
\textbf{\textcolor{maroon}{Tools \& Collaboration:}}\\[4pt]
\cvtag{Power BI}
\cvtag{Jira}
\cvtag{Confluence}
\cvtag{GitHub}
\cvtag{Bitbucket}
\cvtag{Azure Repos}
\cvtag{MS Office Suite}
\cvtag{Outlook}
\cvtag{MS Teams}
\cvtag{SharePoint}
\cvtag{ManageEngine SP}
\cvtag{Zendesk}
\cvtag{Workday}
\cvtag{Copilot}

\cvsection{Languages}

\cvskill{English}{5}
\divider

\cvskill{German}{2} %% Supports X.5 values.
\divider

\cvskill{Hindi}{5}

%% Yeah I didn't spend too much time making all the
%% spacing consistent... sorry. Use \smallskip, \medskip,
%% \bigskip, \vspace etc to make adjustments.
\medskip

% \divider

% \cvsection{Referees}

% \cvref{name}{email}{mailing address}
% \cvref{Mr. Fahad Islam}{NIS Biotech}{salesnisbiotech@gmail.com}
% {\faPhone\ +919045673320}
% {Plot F-11, Industrial Area, UPSIDC, MG Road, Ghaziabad 201015, India}

\end{paracol}


\end{document}
